
\documentclass{article}
\usepackage{colortbl}
\usepackage{makecell}
\usepackage{multirow}
\usepackage{supertabular}

\begin{document}

\newcounter{utterance}

\centering \large Interaction Transcript for game `wordle', experiment `medium\_frequency\_words\_no\_clue\_no\_critic', episode 8 with reflexion{-}Qwen3.5{-}9B.
\vspace{24pt}

{ \footnotesize  \setcounter{utterance}{1}
\setlength{\tabcolsep}{0pt}
\begin{supertabular}{c@{$\;$}|p{.15\linewidth}@{}p{.15\linewidth}p{.15\linewidth}p{.15\linewidth}p{.15\linewidth}p{.15\linewidth}}
    \# & \multicolumn{2}{c}{Player} && \multicolumn{2}{c}{Game Master} \\
    \hline

    \theutterance \stepcounter{utterance}  
    & & & \multicolumn{4}{p{0.6\linewidth}}{
        \cellcolor[rgb]{0.9,0.9,0.9}{
            \makecell[{{p{\linewidth}}}]{
                \texttt{\tiny{[P1$\langle$GM]}}
                \texttt{You have attempted this same task before and did not succeed. These are the plans you wrote after those attempts.} \\
\texttt{Trial \#1: New plan: I failed because I started generating a long, rambling explanation and an incomplete word ("r\_n") instead of providing a concise, valid five{-}letter word in the required format. Next time, I will strictly adhere to the format: provide a brief explanation of the logic derived from the feedback, followed immediately by a complete, valid five{-}letter English word that fits the current constraints (e.g., pattern \_e\_r\_ with known incorrect letters excluded). I will ensure the word is fully spelled out and contains no red{-}letter candidates.} \\
\\ 
\texttt{Trial \#2: New plan: I failed because I got lost in a complex internal monologue about letter positions and generated an incomplete word ("r\_n") instead of a valid five{-}letter guess. Next time, I will strictly adhere to the required format: provide a concise explanation of the logic derived from the feedback, followed immediately by a complete, valid five{-}letter English word that fits the current constraints (e.g., pattern \_e\_r\_ with known incorrect letters excluded). I will ensure the word is fully spelled out and contains no red{-}letter candidates.} \\
\\ 
\texttt{Trial \#3: New plan: I failed because I got stuck in an infinite loop by repeatedly guessing "never" after realizing it was incorrect, instead of generating a new valid word that fits the pattern `\_e\_e\_r` while excluding the known incorrect letters (n, v, r, a, c, l, p). Next time, upon receiving feedback that confirms 'e' is in positions 2 and 4, and 'n', 'v', 'r' are misplaced or incorrect, I will immediately construct a new five{-}letter word using the pattern `\_e\_e\_r` with different starting and ending letters (e.g., "beers", "peers", "seers", "teers", "zeers") to break the loop and find the target. I will ensure the explanation explicitly states the new candidate word chosen to satisfy the pattern constraints.} \\
\\ 
\\ 
\texttt{You are a language wizard who likes to guess words by using the given rules.} \\
\\ 
\texttt{Welcome to Wordle! You have six attempts to guess the target word, a valid English word of five lowercase letters (a{-}z). Please use the tags "explanation:" and "guess:" to provide a concise explanation for each guess.} \\
\\ 
\texttt{For instance, if your guess is "apple", your response should be} \\
\texttt{explanation: this is a common five{-}letter English word, and I am starting my guess with this word.} \\
\texttt{guess: apple} \\
\\ 
\texttt{After each guess, your answer will be validated, and you will receive feedback indicating which letters are correct (green), which letters are correct but in the wrong position (yellow), and which letters are incorrect (red). This feedback can be useful in determining which letters to include or exclude in your next guess.} \\
\\ 
\texttt{For example, the feedback for "apple" might be:} \\
\texttt{guess\_feedback: a<yellow> p<yellow> p<green> l<yellow> e<red>} \\
\\ 
\texttt{The explanation should contain details about how the guess\_feedback is used to arrive at a new guess.} \\
\\ 
\texttt{Do not generate any other text.} \\
\\ 
\texttt{Let's begin with your first guess.} \\
            }
        }
    }
     \\ \\

    \theutterance \stepcounter{utterance}  
    & \multicolumn{4}{p{0.6\linewidth}}{
        \cellcolor[rgb]{0.561,0.737,0.561}{
            \makecell[{{p{\linewidth}}}]{
                \texttt{\tiny{[P1$\rangle$GM]}}
                \texttt{explanation: I will start with a common five{-}letter word that contains a mix of frequent vowels and consonants to maximize information gain in the first guess.} \\
\texttt{guess: apple} \\
            }
        }
    }
    & & \\ \\

    \theutterance \stepcounter{utterance}  
    & & & \multicolumn{2}{p{0.3\linewidth}}{
        \cellcolor[rgb]{0.9,0.9,0.9}{
            \makecell[{{p{\linewidth}}}]{
                \texttt{\tiny{[GM$|$GM]}}
                \texttt{attempts = 1} \\
\texttt{target\_word = renew} \\
\texttt{guess = apple} \\
\texttt{guess\_feedback = a<red> p<red> p<red> l<red> e<yellow>} \\
            }
        }
    }
    & & \\ \\

    \theutterance \stepcounter{utterance}  
    & & & \multicolumn{4}{p{0.6\linewidth}}{
        \cellcolor[rgb]{0.9,0.9,0.9}{
            \makecell[{{p{\linewidth}}}]{
                \texttt{\tiny{[P1$\langle$GM]}}
                \texttt{guess\_feedback: a<red> p<red> p<red> l<red> e<yellow>} \\
\\ 
\texttt{Provide your response only in this format.} \\
\texttt{explanation: details} \\
\texttt{guess: word} \\
            }
        }
    }
     \\ \\

    \theutterance \stepcounter{utterance}  
    & \multicolumn{4}{p{0.6\linewidth}}{
        \cellcolor[rgb]{0.561,0.737,0.561}{
            \makecell[{{p{\linewidth}}}]{
                \texttt{\tiny{[P1$\rangle$GM]}}
                \texttt{explanation: The feedback indicates that 'a', 'p', 'l' are incorrect, while 'e' is correct but in the wrong position (not at the end). I will exclude 'a', 'p', 'l' and try a new word with 'e' in a different position, such as position 2.} \\
\texttt{guess: better} \\
            }
        }
    }
    & & \\ \\

    \theutterance \stepcounter{utterance}  
    & & & \multicolumn{2}{p{0.3\linewidth}}{
        \cellcolor[rgb]{0.9,0.9,0.9}{
            \makecell[{{p{\linewidth}}}]{
                \texttt{\tiny{[GM$|$GM]}}
                \texttt{Error: The length of the guessed word is not 5.} \\
            }
        }
    }
    & & \\ \\

    \theutterance \stepcounter{utterance}  
    & & & \multicolumn{2}{p{0.3\linewidth}}{
        \cellcolor[rgb]{0.9,0.9,0.9}{
            \makecell[{{p{\linewidth}}}]{
                \texttt{\tiny{[GM$|$GM]}}
                \texttt{game\_result = ABORT} \\
            }
        }
    }
    & & \\ \\

    \theutterance \stepcounter{utterance}  
    & & \multicolumn{4}{p{0.6\linewidth}}{
        \cellcolor[rgb]{0.95,0.95,0.95}{
            \makecell[{{p{\linewidth}}}]{
                \texttt{\tiny{[P1$\leftrightarrow$P]}}
                \texttt{Post{-}episode Reflexion (not sent to GM):} \\
\texttt{New plan: I failed because my explanation incorrectly stated that 'l' was an incorrect letter (red), when the feedback `a<red> p<red> p<red> l<red> e<yellow>` clearly shows 'l' is also red. Consequently, I guessed "better," which contains 'l' and violates the constraint to exclude red letters. Next time, I will strictly parse the feedback tags: exclude all letters marked `<red>` (a, p, l) and only keep or reposition letters marked `<yellow>` or `<green>`. I will construct a new five{-}letter word using only the valid letters (e.g., using 'e' in a new position) and ensure the explanation accurately reflects the specific letters to exclude before generating the guess.} \\
            }
        }
    }
    & \\ \\

\end{supertabular}
}

\end{document}
